\section{The binary parity theorem}
\label{sec:parity-proof}

The even base case is immediate.  At $k=2$ the set
\[
  M_2=\{00,01,11\}
\]
selects one word from each of the three PCRs, is RC invariant, and leaves
only the vertex $10$, with no residual loop.  It is therefore an ordinary
MDS.

\begin{proof}[Proof of Theorem~\ref{thm:parity}]
For even $k\geq4$, Theorem~\ref{thm:even} supplies an RC-invariant ordinary
MDS; the preceding paragraph covers $k=2$.  For odd $k\geq3$,
Theorem~\ref{thm:odd}, specialized to the binary complement, rules out an
RC-invariant decycling set at the ordinary minimum.  These cases exhaust the
stated domain.
\end{proof}

The logical dependencies are worth making explicit:
\[
\begin{array}{c}
\text{ordinary PCR lower bound and MDS construction}\quad+
\quad\text{valid $F$-move preservation}
\\[2mm]
\Downarrow
\\[-1mm]
\text{arbitrary-zero-tie seed}
\;\Longrightarrow\;
\text{exact half-plane coboundary}
\;\Longrightarrow\;
\text{finite legal defect descent}
\\[2mm]
\Downarrow
\\[-1mm]
\text{all even binary orders, combined independently with the odd bridge.}
\end{array}
\]
No component-count conjecture, solver outcome, bounded-order pattern, or
connectivity claim enters this chain.

It is useful to name the quantity separated by the theorem.  Let
$\tau_{\RC}(k)$ denote the minimum cardinality of an RC-invariant decycling
set in $B(2,k)$ and define the \emph{symmetry premium}
\[
  \Delta_2(k)=\tau_{\RC}(k)-N_2(k).
\]
The theorem says exactly that $\Delta_2(k)=0$ for even $k$, whereas
Theorem~\ref{thm:odd} gives $\Delta_2(k)\geq2$ for odd $k$.  No exact formula
for the odd premium is asserted.
